%% ============================================================================
%%  05_modelo_datos/contenido.tex
%%
%%  Subdocumento 5 del Formulario T-7 — Modelo y gestión de datos.
%%
%%  Estructura:
%%    5.1  Modelo — Dominios de información y modelo de datos por dominio.
%%    5.2  Gestión de datos — Motor, paradigma, separación OLTP/OLAP, calidad.
%%    5.3  Estrategia de migración — Saneamiento, validación y conciliación.
%%    5.4  Estrategia de desempeño — Indexación, particionamiento, caché.
%%
%%  Las secciones complementarias (trazabilidad, GS1/EDI, diccionario y
%%  evidencia de entrega) se conservan en el orden original como 5.5 a 5.8.
%%
%%  Tablas: 05_modelo_datos/tablas/*.tex (fragmentos traídos con \parte{}).
%% ============================================================================

%  El número de capítulo se ancla al número de subdocumento del T-7.
\setcounter{chapter}{4}
\chapter{Modelo y gestión de datos}
\label{cap:datos}

% ── Texto de introducción ────────────────────────────────────────────────────

Este capítulo, define cómo se organiza, almacena, migra y optimiza la
información que la operación de Puelche produce, consume y conserva.  Su
alcance cubre desde la volumetría de partida ---14.200 clientes, 31.000
pedidos mensuales con 260.000 líneas y 2,4 millones de unidades, 1.400
entregas diarias con un pico de septiembre cercano a 2.600--- hasta la
política de eliminación segura al término del contrato.

La arquitectura de datos se describe conforme a ISO/IEC/IEEE 42010, en
coherencia con la arquitectura lógica y con la arquitectura física del
Subdocumento~4, y responde a los requisitos del capítulo 15 del caso y a las
Bases Técnicas Transversales.

\begin{description}
\item[Sección 5.1 Modelo.] Presenta los dominios de información, los entornos de
  datos y las quince bases lógicas que sostienen la operación, con su relación
  con los doce módulos funcionales de la arquitectura lógica.

\item[Sección 5.2 Gestión de datos.] Justifica la selección de motor y paradigma de
  persistencia por dominio, conforme al teorema CAP, declara la separación
  entre almacenamiento transaccional y analítico (RT-05.05), y establece las
  políticas de calidad, retención, archivado y eliminación segura.

\item[Sección 5.3 Estrategia de migración.] Desarrolla la migración, saneamiento,
  validación y conciliación de los datos históricos procedentes del ERP de 2017,
  el sistema de almacenes de 2013 y las planillas, con volúmenes y ventanas de
  corte compatibles con el cronograma del Artículo 17.

\item[Sección 5.4 Estrategia de desempeño.] Fundamenta la indexación,
  particionamiento, caché y optimización de consultas en la volumetría del
  caso, con foco en la ventana crítica de despacho de 05:30 a 07:00.
\end{description}

El diccionario de datos detallado ---con nombre, tipo, dominio de valores,
obligatoriedad, dueño y sensibilidad por atributo--- se entrega como anexo
(RT-05.01).  Las secciones complementarias de trazabilidad sanitaria
(sección 5.5), estándares GS1 e intercambio electrónico (sección 5.6), diccionario y
catálogo (sección 5.7), y evidencia de entrega digital (sección 5.8) completan el
capítulo.

\nota{La política de retención se declara contra RT-16.10. El capítulo 15 del
caso la rotula RT-05.10, que en las Bases Técnicas Transversales corresponde a
catálogo de datos y linaje y es deseable, no obligatorio. El desajuste se
responde en el Formulario T-12 contra el código correcto y se eleva como
consulta conforme al Artículo 43.3. Ver sección 5.2.3.}


% ============================================================================
%  5.1  MODELO — DOMINIOS DE INFORMACIÓN
% ============================================================================
\section{Modelo: dominios de información y modelo de datos por dominio}
\label{sec:modelo-dominios}

El dominio de información cubre la totalidad de los datos que produce la
operación: preventa y pedidos, inventario y bodega en dos centros de
distribución y tres plataformas de cross-docking, preparación, reparto y
entregas, cobranza y cartera, trazabilidad sanitaria y cadena de frío,
telemetría de flota, canal moderno, documentos tributarios, analítica de
gestión y datos maestros.

La volumetría de referencia proviene de la Tabla 14.1 y del capítulo 15 del
caso: 14.200 clientes; 31.000 pedidos al mes, con 260.000 líneas y 2,4
millones de unidades; 180 proveedores; unos 34.000 documentos tributarios
mensuales; 9.500 SKU proyectados a tres años; telemetría de 28 sensores y 18
termógrafos; y del orden de 1.400 entregas diarias, con un peak de septiembre
cercano a 2.600.


% ── 5.1.1  Entornos de datos ─────────────────────────────────────────────────
\subsection{Entornos de datos}
\label{sec:entornos-datos}

\parte{tablas/entornos}

La topología es multisitio por parametrización (RF-02.02, RT-02.12): cada
centro informa su stock al consolidado y un sitio nuevo se incorpora
replicando la plantilla del clúster, sin rediseñar el modelo de datos.  El
emplazamiento híbrido ---on-premise para bodega y nube para preventa, reparto
y analítica--- se describe en la arquitectura física del Subdocumento~4; la
separación de roles lógicos por sitio se detalla en la Capa~6 de la
arquitectura lógica.

% ── 5.1.2  Modelo de datos conceptual por dominio ────────────────────────────
\subsection{Modelo de datos conceptual por dominio}
\label{sec:modelo-conceptual}

El modelo de datos se fundamenta en las entidades de negocio canónicas
definidas en la arquitectura lógica, organizadas en quince bases lógicas.
Las entidades principales y sus eventos se resumen a continuación; el
diccionario completo por atributo se entrega como anexo (RT-05.01).

\begin{description}
\item[Cliente.] RUT, razón social, canal (tradicional, food service, cadenas),
  crédito, georreferencia y estado de bloqueo.  Eventos: creación,
  actualización y cambio de bloqueo.

\item[Pedido.] Identificador UUID, preventista, cliente, líneas de detalle
  (SKU, cantidad, precio) y ciclo de vida (tomado $\to$ confirmado $\to$
  preparado $\to$ despachado $\to$ entregado $\to$ rendido).  Eventos: toma,
  confirmación, asignación de línea, despacho y evidencia de entrega.

\item[SKU / Producto.] Codificación GTIN/GS1, unidad, clasificación de
  temperatura, rangos térmicos (RF-09.01) y vida útil.

\item[Lote.] Unidad de trazabilidad sanitaria, vinculada al lote del
  proveedor (GS1: GTIN + lote + vencimiento) y enlazada a la unidad logística
  SSCC en cada movimiento interno.  Eventos: recepción, bloqueo, inicio de
  retiro sanitario y enlace con unidad logística.

\item[Misión de preparación.] HHT, oleada, ubicación, unidades y ciclo de vida
  (descargada $\to$ ejecutada $\to$ cerrada).

\item[Ruta / viaje.] Planificador, vehículo, conductor, ventanas, secuencia de
  entregas y geocercas.

\item[Entrega.] Pedido, viaje, prueba de entrega (POD: firma, QR, fotografía),
  documentos DTE y efectivo.  Eventos: evidencia, reintento y aprobación de
  rendición.

\item[Documento tributario.] Factura, boleta, guía de despacho, folio del SII
  y acuse.

\item[Envase retornable.] Canastillo o pallet, cliente, saldo y pérdida
  estimada del 14\,\% anual (decisión 16.1 N.\textdegree\,10).  Control por
  cuenta corriente por cliente, no por unidad identificada.

\item[Sensor / registro térmico.] Dispositivo, lote o posición, temperatura y
  excursión térmica.
\end{description}

Los eventos usan verbos en pasado, son inmutables y constituyen la base de los
esquemas AsyncAPI 2.6 de la capa de integración.  Toda escritura sin conexión
es idempotente por UUID (RT-02.06).

El almacenamiento se distribuye en quince dominios lógicos, cada uno con un
propietario funcional trazable a los módulos M1--M12 de la arquitectura
lógica.  Esta separación incluye bases transaccionales, almacenes analíticos,
caché y objetos; no representa quince servidores físicos ni quince motores de
base de datos independientes.


% ============================================================================
%  5.2  GESTIÓN DE DATOS
% ============================================================================
\section{Gestión de datos}
\label{sec:gestion-datos}

% ── 5.2.1  Motor y paradigma de persistencia por dominio ─────────────────────
\subsection{Motor y paradigma de persistencia por dominio}
\label{sec:motor-paradigma}

La justificación de paradigma y motor responde a RT-05.02.  El punto de
partida es que la operación impone cuatro exigencias incompatibles entre sí,
cada una con una posición distinta en el triángulo del teorema CAP:

\begin{enumerate}
\item \textbf{Transaccional de bodega y preventa} (stock, crédito, reservas):
  exige consistencia fuerte y reconciliación determinista después de un corte.
  Posición CAP: \emph{CP} ---consistencia y tolerancia a partición--- con
  disponibilidad local garantizada por autonomía de 24 horas.

\item \textbf{Ingesta de sensores y termógrafos} de cadena de frío: escritura
  masiva de baja latencia sobre datos volátiles que no deben saturar el motor
  transaccional (RNF-11.02).  Posición CAP: \emph{AP} ---disponibilidad y
  tolerancia a partición, con consistencia eventual.

\item \textbf{Serie de tiempo consolidada} de temperatura y telemetría:
  analítica columnar con retención de cinco años y consulta por rango.
  Posición CAP: \emph{AP} columnar.

\item \textbf{Evidencia documental} ---prueba de entrega y documentos
  tributarios---: inmutabilidad legal durante seis años, con bloqueo de objeto
  en modo cumplimiento (RT-16.07).  Posición: \emph{CP} con disponibilidad
  provista por replicación de objetos.
\end{enumerate}

Ningún motor único cubre las cuatro posiciones del teorema CAP sin sacrificar
desempeño (RT-05.02).  Por eso la decisión de arquitectura ADR-04 establece
persistencia políglota por dominio, declarando motor y posición CAP para cada
clase de dato.  La decisión general es disponibilidad con consistencia eventual
reconciliada, y consistencia fuerte únicamente donde el negocio la exige.

\parte{tablas/persistencia}

La elección de motor por dominio corresponde a la decisión 16.1
N.\textdegree\,18.  El volumen total es manejable en PostgreSQL con índices y
PostGIS para georreferencia; el motor clave-valor absorbe el pico de ingesta
IoT; y la serie consolidada vive en la capa analítica, sin añadir un motor de
series dedicado.  Se descartó InfluxDB precisamente por eso: sería un segundo
motor que operar con un equipo de TI de cuatro personas.  Las alternativas
evaluadas y su fundamento constan en ADR-04.

El modelo materializa además el despliegue híbrido del Artículo
16\textsuperscript{o}.  El maestro de bodega vive on-premise, con PostgreSQL y
PostGIS por sitio y autonomía de 24 horas sin enlace (RT-03.10), y se replica
por captura de cambios sobre el registro de escritura anticipada hacia la
instancia de nube, que concentra el transaccional de preventa, reparto, canal
moderno e inteligencia de negocio.  El terreno opera sin red de primera clase:
las aplicaciones mantienen una fotografía local de los datos de lectura
---stock, precios, crédito y rutas--- y un búfer de escrituras idempotente que
se concilia al reconectar (RT-03.12), sin perder ni duplicar pedidos.  El
detalle del dimensionamiento de la infraestructura de cada motor se documenta
en la arquitectura física del Subdocumento~4.


% ── 5.2.2  Separación transaccional y analítica ─────────────────────────────
\subsection{Separación entre almacenamiento transaccional y analítico}
\label{sec:separacion-oltp-olap}

RT-05.05 es obligatorio: el almacenamiento transaccional y el analítico están
separados por diseño, y ninguna consulta analítica puede degradar la operación.

\parte{tablas/separacion}

El flujo de datos hacia la capa analítica es unidireccional: los eventos
incrementales de la capa de integración alimentan el almacén analítico
(Redshift Serverless sobre S3 Parquet, procesado con AWS Glue), y los procesos
de flujo pesado se ejecutan fuera de la ventana crítica de 05:30 a 07:00.
Los tableros del módulo de inteligencia de negocio (M10) consultan
exclusivamente la capa analítica, nunca el transaccional en vivo.

El modelo dimensional se estructura en hechos ---ventas, entregas, stock y
costo de servir--- y dimensiones ---cliente, SKU, tiempo, ruta y canal--- en
esquema estrella sobre la capa columnar, con descenso hasta la línea de pedido
y el lote para la trazabilidad sanitaria completa (RT-05.29).

Las latencias comprometidas ---operación del día en cinco minutos o menos,
cierre comercial en dos horas y gestión en cuatro--- se verifican como
objetivos de nivel de servicio en la capa de observabilidad y se prueban con
carga de ventana durante la marcha blanca.


% ── 5.2.3  Calidad, retención, archivado y eliminación ──────────────────────
\subsection{Calidad de datos, retención, archivado y eliminación segura}
\label{sec:calidad-retencion}

\subsubsection{Calidad de datos}

\parte{tablas/calidad}

Las reglas de calidad se declaran en el diccionario (RT-05.01) y se monitorean
con controles automáticos de duplicados, saltos de secuencia y totales de
conciliación, conforme al marco ISO/IEC 25012.  El incumplimiento genera una
tarea en el flujo de excepciones y el tablero de calidad queda disponible para
el CLIENTE (RT-05.04).

\subsubsection{Retención, archivado y eliminación}

\parte{tablas/retencion}

La eliminación segura (RT-05.07) sigue un procedimiento verificable al
vencimiento del plazo, con registro inalterable (RT-16.07), borrado seguro de
los medios que salen de servicio y disposición final a través de un gestor
autorizado.

La exportabilidad (RT-05.06) garantiza que la totalidad de la información del
CLIENTE se exporte en formatos abiertos y documentados ---CSV, JSON y Parquet
con su diccionario---, en cualquier momento del contrato, sin costo y sin
intervención del proponente.  Al término del contrato, el Artículo
85\textsuperscript{o} y RT-05.08 obligan a entregar copia íntegra en formato
interoperable con diccionario y a eliminar de forma certificada los datos del
mandante, con el directorio de tratamiento actualizado y auditoría.


% ============================================================================
%  5.3  ESTRATEGIA DE MIGRACIÓN
% ============================================================================
\section{Estrategia de migración}
\label{sec:migracion}

La migración cubre los datos históricos de tres fuentes legadas: el ERP de
2017, el sistema de almacenes de 2013 y las planillas de operación.  Los
volúmenes y la retención posterior se resumen en la tabla siguiente.

\parte{tablas/migracion}

\subsection{Metodología de migración}
\label{sec:migracion-metodologia}

\begin{description}

\item[Perfilado y saneamiento previo (RT-05.12).] Etapa de perfilado sobre los
datos de origen con informe de defectos y decisión sobre cada uno: corregir,
aceptar con registro o excluir.  Los casos esperados están declarados: clientes
duplicados por RUT, campos de lote en texto libre o vacíos ---el retiro de
marzo mostró un 41~\% de recepciones sin lote---, equivalencias de maestro
entre cadenas y saldos de envases sin control previo.

\item[Reglas de transformación y carga verificada.] Los datos maestros se
cargan con validación en el punto de captura y conciliación por recuentos.  La
carga transaccional inicial se hace por olas y por sitio, con estrategia azul y
verde y plan de reversión.

\item[Dos ensayos completos sobre preproducción (RT-05.13).] Con medición del
tiempo total y del resultado de la conciliación: un ensayo y un reensayo tras
la corrección.

\item[Conciliación cuantitativa (RT-05.14).] Recuentos, sumas de control y
muestreo dirigido.  Toda diferencia queda explicada y registrada en la bitácora
de reconciliación.  La reposición del conteo cíclico apuntala la exactitud del
stock migrado, con meta de bajar la diferencia del 2,3~\% actual a menos del
1~\%.

\item[Históricos no migrados (RT-05.15).] Quedan accesibles en un repositorio
de consulta de solo lectura durante el período de retención, sin depender del
sistema legado, y exportables bajo el régimen del Artículo 85\textsuperscript{o}
en formato interoperable con su diccionario.

\item[Coexistencia con el sistema de almacenes de 2013.] El legado queda detrás
de la misma capa anticorrupción (ADR-11, A-04), con una tabla de capacidad
absorbida ---recepción GS1, asignación de ubicación, misiones con terminal
portátil y conteo cíclico--- que se vacía ola a ola por sitio, con plan de
reversión y retiro al cierre de la Etapa 1.

\end{description}

\subsection{Ventanas de corte y compatibilidad con el cronograma}
\label{sec:migracion-ventanas}

La migración se ejecuta conforme a las etapas del Artículo 17: la Etapa~1
comprende los meses 1--15, con producción en el mes 16; la Etapa~2, los
meses 13--20, con producción en el mes 21.  Las cargas masivas se programan
fuera de la ventana crítica de 05:30 a 07:00, con estrategia azul-verde por
sitio que permite la reversión sin pérdida de operación.  Los dos ensayos de
preproducción deben completarse antes del corte y su resultado ---tiempo total
y tasa de conciliación--- se documenta como criterio de aceptación del paso a
producción.

La coexistencia con los legados durante la transición se gobierna por la capa
anticorrupción, que traduce los contratos de la solución al formato del
legado sin reemplazarlo ni modificar su código.  La tabla de capacidad
absorbida se vacía ola a ola por sitio, de modo que el riesgo de la migración
se acota a un dominio funcional y a una instalación a la vez.


% ============================================================================
%  5.4  ESTRATEGIA DE DESEMPEÑO
% ============================================================================
\section{Estrategia de desempeño}
\label{sec:desempeno}

La estrategia de desempeño se fundamenta en la volumetría del caso: 31.000
pedidos mensuales, 260.000 líneas, 2,4 millones de unidades, 1.400 entregas
diarias ---con un peak de septiembre de 2.600--- y un dimensionamiento del
orden de 105 transacciones por segundo de base, con margen de crecimiento a
tres años (RT-09.03).  El diseño se dimensiona contra el pico de septiembre en
la ventana de 05:30 a 07:00, nunca contra el promedio.

\subsection{Indexación}
\label{sec:indexacion}

Las tablas transaccionales llevan claves foráneas completas.  Existen índices
por lote GS1 ---identificador de producto, lote y vencimiento---, por unidad
logística y por identificador de despacho: es la estructura que sostiene el
retiro sanitario en menos de dos horas.  Las geometrías, rutas y geocercas
usan índices espaciales GiST.  La trazabilidad hacia adelante y hacia atrás se
apoya en índices compuestos sobre la base de calidad y trazabilidad.

Las réplicas de solo lectura atienden las consultas de negocio sin competir con
la escritura del maestro (RT-09.05).  En la capa analítica, la distribución se
hace por clave de negocio ---ruta o lote--- con compresión columnar sobre los
hechos de ventas, entregas, stock y costo de servir, y sus dimensiones de
cliente, SKU, tiempo, ruta y canal.

\subsection{Particionamiento}
\label{sec:particionamiento}

Las transacciones de bodega ---recepción, misiones y despacho--- se particionan
por mes en el PostgreSQL on-premise.  Es el primer remedio declarado ante el
cuello de botella de escritura de la ventana 05:30--07:00, dentro de una
progresión explícita: escala vertical, luego particionado con agrupador de
conexiones, luego réplica de lectura y finalmente un cuarto nodo o
fragmentación.

La mensajería garantiza orden por partición, con clave de agrupación por sitio
o entidad, lo que permite reconciliar entre sitios y ordenar los eventos de
trazabilidad (RT-02.07).  La sincronización de terreno usa partición y
reanudación por bloques para cumplir los diez minutos de turno y las dos horas
del centro de distribución tras reconectar.

La telemetría en bruto se retiene 30 días por vencimiento automático de
partición; la serie consolidada se particiona por mes en almacenamiento
columnar, con cinco años de retención.  El concentrador de intercambio
electrónico particiona por cadena, mediante perfiles de configuración y no
mediante desarrollo.

\subsection{Caché}
\label{sec:cache}

La caché en memoria (Redis sobre Amazon ElastiCache) mantiene stock caliente,
precios y sesiones de inicio único para sostener la consulta de stock y crédito
bajo dos segundos en preventa (RT-09.01), con patrón de lectura diferida y
regeneración transparente en menos de quince minutos ante una pérdida.

En el dispositivo, la caché de turno se precarga al inicio ---saldo, stock y
tarifa, RT-17.03--- y sostiene el turno completo sin señal dentro de un consumo
objetivo de entre 8 y 12 MB.  La identidad se cachea localmente con vigencia de
8 horas en bodega y 14 horas en reparto, para que la autenticación no dependa
de la red.  Los tableros se cachean en las consolas locales de contingencia
(RT-03.13).

\subsection{Optimización operacional}
\label{sec:optimizacion}

El agrupador de conexiones (PgBouncer) y la pregeneración nocturna de
documentos de despacho y de validaciones por lote de crédito y stock descargan
la ventana crítica.  Los umbrales de RT-09.01 se verifican bajo carga de
ventana durante la marcha blanca: confirmación de línea de picking en un
segundo o menos, registro de entrega en dos segundos, línea de preventa en
1,5 segundos y consulta de stock o crédito en dos segundos, todos en percentil
95, frente al dimensionamiento declarado.

La progresión de escalado está documentada: escala vertical del nodo PostgreSQL,
luego particionado con agrupador de conexiones, luego réplica de lectura y
finalmente fragmentación.  Cada escalón tiene un umbral de activación medido
en la capa de observabilidad, de modo que el equipo TI de cuatro personas
dispone de un camino predecible para absorber el crecimiento.


% ============================================================================
%  SECCIONES COMPLEMENTARIAS (5.5 – 5.8)
% ============================================================================

\section{Trazabilidad sanitaria y retiro en menos de dos horas}
\label{sec:trazabilidad}

\subsection{Unidad de trazabilidad}

La identidad primaria de la trazabilidad es el lote del proveedor en estándar
GS1: identificador de producto, número de lote y fecha de vencimiento, con los
rangos térmicos asociados.  La unidad logística se enlaza al lote en cada
movimiento interno ---recepción, ubicación, picking y unidad de despacho---.

Se descartaron la caja y el pallet como identidad primaria por dos razones.  La
primera es de volumen: 9.500 SKU y 2,4 millones de unidades al mes.  La segunda
es de diagnóstico: el 41~\% de las recepciones del producto retirado no tenía
lote, de modo que el problema es la captura y no el nivel de agregación.  Por
eso la recepción exige la lectura del lote del proveedor, el 100~\% de las
recepciones de productos que lo requieren registra el lote, y la etiqueta de
unidad logística se imprime al armar el pallet.

\subsection{Modelo de eventos}

La trazabilidad se registra evento a evento conforme a GS1 EPCIS (RT-05.23).
Cada movimiento ---recibido, ubicado, preparado, despachado, entregado---
genera un evento idempotente con actor, equipo o dispositivo, marca de tiempo y
valores anteriores y posteriores (RT-05.03), conservado como registro inmutable
(RT-16.07).

La bidireccionalidad es completa.  Hacia adelante, desde un lote del proveedor
se responde con evidencia, y no con estimación, a qué clientes llegó, en qué
fecha, en qué cantidad y con qué documento tributario.  Hacia atrás, desde una
unidad en el local del cliente se llega a la recepción que la originó.  La
analítica permite descender hasta la línea de pedido y el lote (RT-05.29).

El retiro sanitario en menos de dos horas ---frente a los nueve días con
resultado estimado de la situación actual--- se resuelve recorriendo la base de
calidad y trazabilidad con índices por lote y unidad logística sobre la fuente
de verdad consolidada, con objetivo de nivel de servicio comprometido y
evidencia exportable ante la autoridad.  A ello se suma el registro continuo de
temperatura en cámaras y vehículos, con detección de excursión en el borde y
regla escrita de bloqueo.

\subsection{Gobernanza}

Los eventos se publican por la capa de integración, con orden por partición y
deduplicación, y alimentan la lectura consolidada en el servidor.  La bitácora
de auditoría registra toda operación con quién, qué, cuándo, con qué
dispositivo y con qué valores previos y posteriores, lo que sostiene tanto la
trazabilidad del dato de negocio como la del lote.  La matriz de trazabilidad
del capítulo 17.1 conserva la relación entre origen, requerimiento, componente
y prueba de cada dato.

\section{Estándares GS1, intercambio electrónico y formato tributario}
\label{sec:estandares}

\parte{tablas/estandares}

La integración sigue el principio de asíncrono por defecto y síncrono por
excepción.  Lo síncrono se restringe a los actos que lo exigen ---emisión
tributaria, autorización de pago y lecturas críticas de stock o crédito---,
siempre con tiempo de espera explícito; todo lo demás circula por colas con
cola de mensajes fallidos y reintentos.  Los perfiles y contratos por
contraparte se versionan y se prueban contra el banco de pruebas de cada una.

\section{Diccionario de datos y catálogo de activos}
\label{sec:diccionario}

RT-05.01 es obligatorio: el diccionario se documenta por base lógica con
formato estándar.

\parte{tablas/diccionario}

El diccionario se publica en el catálogo de datos con versionado, alimenta la
matriz de trazabilidad del capítulo 17.1, habilita el linaje sobre la capa
analítica y acompaña la copia íntegra interoperable que exige el Artículo
85\textsuperscript{o} al término del contrato.  El detalle de cada atributo
---nombre, tipo, dominio de valores, obligatoriedad, propietario y
sensibilidad--- queda cerrado en el entregable formal del diccionario, previo a
producción.

\section{Evidencia de entrega digital y acuse de recibo}
\label{sec:evidencia}

Esta sección responde al criterio 8 del capítulo 18 ---prueba de entrega
digital disponible el mismo día, frente a los doce días que hoy toma resolver
un reclamo con guía de papel---, a RT-16.14 ---el receptor habitual es una
persona natural sin firma electrónica avanzada--- y a la decisión de diseño
16.1 N.\textdegree\,15, que pregunta cómo se prueba la entrega cuando quien
recibe no es el titular, no sabe firmar en pantalla o se niega a hacerlo.

\parte{tablas/evidencia}

La evidencia capturada ---firma, fotografías y código QR--- se conserva con
bloqueo de objeto y retención de seis años, de forma inmutable (RT-16.07), y se
articula con la guía de despacho electrónica y su acuse de recibo.  El acuse de
recibo de las mercaderías es el documento que da mérito ejecutivo a la factura,
de modo que la solución no puede comprometer sus efectos legales.

La firma se aplica conforme a la Ley 19.799 y a RT-16.17 y RT-16.18, con
verificación de validez del certificado al momento de firmar, sello de tiempo y
evidencia verificable después del vencimiento.  La evidencia local se captura
sin conexión, el documento se timbra al sincronizar con folio reservado por
dispositivo (RT-03.12) y el acuse se envía al cliente el mismo día, con
confirmación por al menos tres canales (RT-16.08).  En el canal moderno el acuse
es digital, dentro de la ventana acordada.
